\begin{papercard}{oloughlin2019}
\cf{Stages} \stagetag{A} \stagetag{D} \stagetag{C}. Transparency review of the privacy/security policies of depression apps. It is notable for a two-stage screening with reported screening agreement and for treating \emph{policy absence} as part of the transparency score.

\cf{Aim} Assess the availability and thoroughness of privacy policies of depression apps, and whether they differ by store or by the type of data collected \pc{oloughlin2019}{3}{The present study sought to un}. Hypothesis: many apps would have no policy or a poor one.

\cf{Search} Single term ``depression'' in iTunes and Google Play on 15 October 2017 in Chicago, following prior mental-health app reviews \pc{oloughlin2019}{3}{searched for apps using the search term}. Google Play shows at most 250 results, so iTunes was also capped at the first 250 to avoid platform bias, giving 500 candidates. The cap is justified by users rarely looking beyond the top 10 \pc{oloughlin2019}{3}{As Google Play presents a}.

\cf{Screening} Three raters (the first three authors) removed duplicates ($n=29$) \pc{oloughlin2019}{3}{and eliminated any apps that were duplicates}. Stage 1 screened store descriptions; stage 2 downloaded and reviewed the app, because some criteria are only observable in the app. Inclusion: support/treatment for depression, English, for adults, collects data. Exclusion: professionals only, not standalone, unavailable, not functioning \pc{oloughlin2019}{3}{Apps met inclusion criteria if they (1) aim to}. Each app was screened by two raters at each stage. Initial agreement was 86.7\,\% (stage 1) and 89.5\,\% (stage 2), and disagreements were discussed to full agreement \pc{oloughlin2019}{3}{Raters obtained an 86.7}. 116 apps were eligible.

\cf{Instrument} A yes/no checklist adapted from the Enlight evaluation tool, with items selected for privacy and basic security and adapted using the APA App Evaluation Model \pc{oloughlin2019}{4}{veloped by adapting questions from Baumel}. ``Yes'' required that the policy \emph{explicitly} states the content; absent information counted as ``No'' because a user would not know it \pc{oloughlin2019}{4}{privacy policy explicitly state the content of the question}. Items (Fig.~2): has a policy; PIN/login; encryption \emph{or} local storage; storage/sharing described; deletion; editing; use without identifiable information \pc{oloughlin2019}{3}{Checklist used to evaluate presence and comprehensiveness}. Scoring maps answers to three levels, acceptable / questionable / unacceptable, using a white/light-grey/dark-grey response grid that distinguishes apps collecting identifiable from non-identifiable information \pc{oloughlin2019}{4}{resulting scores were grouped into three categories}. No technical audit was done, so only the comprehensiveness of the \emph{documentation} was assessed \pc{oloughlin2019}{4}{we did not conduct a technical audit}.

\cf{Reliability} The first 23 apps were rated by two raters. ICCs of 0.923--1.00 were judged ``excellent'', and the remaining apps were rated by one rater \pc{oloughlin2019}{4}{Consistency of ratings between raters was con}.

\cf{Key findings} 4\,\% acceptable, 28\,\% questionable, 68\,\% unacceptable, mainly because only 49\,\% had a policy at all \pc{oloughlin2019}{4}{Of the 116 eligible apps, 4}. Policies were found most often in the store (79\,\%) and least often in the app (53\,\%). In-app policies were rarely shown before data collection (11\,\%) \pc{oloughlin2019}{4}{Privacy policies provided in the app were rarely}. iOS apps were more likely to have a policy, $\chi^2(1)=6.07$ \pc{oloughlin2019}{4}{The availability of privacy policies differed significantly by platform}. Apps collecting identifiable information more often had a policy (79\,\% vs.\ 34\,\%).

\cf{Limitations stated} Apps that appeared not to collect data were excluded, although they may collect background data. Policies were not audited against practice, so quality may be overestimated. Comprehensibility was not assessed \pc{oloughlin2019}{5}{We excluded apps that did not}.

\cf{Use for this thesis} Directly relevant to the meeting decision on \emph{apps without policies}: here absence is part of the outcome, so the denominator stays at all eligible apps. Reusable: the explicit ``Yes only if explicitly stated'' rule, the two-stage screening with reported agreement, and the equal result caps per store. Caution: rating only a first block of 23 apps for reliability and then switching to a single rater is a weak design for the thesis, where agreement should be measured on a frozen manual and a planned sample.
\end{papercard}
